% Source: https://artofproblemsolving.com/community/c3h2474166_can_you_manipulate_countable_infinity
% Stable identity: aops-topic-2474166-q01
\textbf{(a)} Let $g(n)=\dfrac n2$ for even $n$ and $g(n)=-\dfrac{n-1}{2}$ for odd $n$, so $g$ lists $\mathbb Z$ as $0,1,-1,2,-2,\dots$. It is a bijection, with inverse $z\mapsto2z$ for $z>0$ and $z\mapsto1-2z$ for $z\le0$ (check: $g(2z)=z$, $g(1-2z)=z$, and these inverse values cover each $n\in\mathbb N$ exactly once).

\textbf{(b)} Every positive integer $N$ can be written as $N=2^{k}q$ with $k\ge0$ and $q$ odd: divide by $2$ as long as possible. This form is unique, because $k$ must be the exponent of $2$ in the prime factorisation of $N$, and then $q=\dfrac{N}{2^k}$. Writing $k=m-1$ and $q=2n-1$ with positive integers $m,n$, each $N$ equals $f(m,n)$ for exactly one pair $(m,n)$. So $f$ is a bijection. $\blacksquare$

\textbf{(c)}
\begin{itemize}
\item $A=\mathbb N$: then $\mathbb N\setminus A=\varnothing$.
\item $A=\{6,7,8,\dots\}$, which is countably infinite via $n\mapsto n+5$: then $\mathbb N\setminus A=\{1,2,3,4,5\}$ has $5$ elements.
\item $A=\{2,4,6,\dots\}$, countably infinite via $n\mapsto2n$: then $\mathbb N\setminus A=\{1,3,5,\dots\}$ is countably infinite via $n\mapsto2n-1$.
\end{itemize}
In each case a set of size $\aleph_0$ is removed from a set of size $\aleph_0$, but the size of what is left can be $0$, $5$ or $\aleph_0$. So subtraction of infinite cardinals has no well-defined answer (and division fails for the same reason), while addition and multiplication behave as in (a) and (b).
